% ==============================================================================
%  3.1  Inferencias sobre los coeficientes de regresión
% ==============================================================================
\section{Inferencias sobre los coeficientes de regresión}
\label{sec:t3-coeficientes}

En todo el tema se supone que es aplicable el modelo de regresión lineal simple con errores
normales (\cref{def:rls}):
\[
  Y_i = \beta_0 + \beta_1X_i + \varepsilon_i, \qquad
  \varepsilon_i \iid \Normal(0,\sigma^2), \quad i=1,\ldots,n .
\]
Hacer inferencias sobre $\beta_0$ y $\beta_1$ consiste en cuantificar la incertidumbre de sus
estimaciones mediante intervalos de confianza y en resolver contrastes de hipótesis sobre sus
valores. Los contrastes de mayor interés son los de la forma $H_0\colon\beta_i=0$, y en
particular
\[
  H_0\colon \beta_1 = 0 \qquad\text{frente a}\qquad H_1\colon \beta_1\neq 0 .
\]
Si $\beta_1=0$, entonces $\E(Y)=\beta_0+0\cdot X=\beta_0$: la recta de regresión es
horizontal y la media de $Y$ no varía con $X$. Más aún, bajo el modelo las $Y_i$ son
independientes con $Y_i\sim\Normal(\beta_0+\beta_1X_i,\sigma^2)$ (\cref{prop:t2-dist-Y}); si
$\beta_1=0$, todas siguen la misma distribución $\Normal(\beta_0,\sigma^2)$, que no depende de
$X$, de modo que \textbf{no hay relación de ningún tipo entre $X$ e $Y$}.

Esta conclusión presupone que el modelo es correcto. Si falla la hipótesis de linealidad, $Y$
puede depender de $X$ de forma no lineal (por ejemplo, cuadrática) aunque la recta que mejor se
ajusta a los datos tenga pendiente nula; por eso es necesario validar el modelo
(\cref{sec:t4-objetivos}).

Para resolver estos contrastes y construir los intervalos es necesario conocer la distribución
en el muestreo de los estimadores $\est{\beta}_0$ y $\est{\beta}_1$, es decir, la distribución de
los valores que toman en muestras repetidas con los valores de $X$ fijos.

\subsection{Distribución en el muestreo de \texorpdfstring{$\est{\beta}_1$}{beta1}}

\begin{proposicion}[$\est{\beta}_1$ como combinación lineal de las $Y_i$][prop:t3-ki]
  El estimador $\est{\beta}_1 = SS_{XY}/\SSX$ puede escribirse como
  \[
    \est{\beta}_1 = \sumi k_iY_i,
    \qquad
    k_i = \frac{X_i-\media{X}}{\sumi (X_i-\media{X})^2} = \frac{X_i-\media{X}}{\SSX},
  \]
  donde los $k_i$ son constantes que dependen únicamente de las $X_i$. Además:
  \[
    \text{(i) } \sumi k_i = 0, \qquad
    \text{(ii) } \sumi k_iX_i = 1, \qquad
    \text{(iii) } \sumi k_i^2 = \frac{1}{\SSX}.
  \]
\end{proposicion}

\begin{demostracion}
  \emph{Expresión de $\est{\beta}_1$.} La suma de las desviaciones respecto de la media es nula:
  \[
    \sumi (X_i-\media{X}) = \sumi X_i - n\media{X} = n\media{X}-n\media{X} = 0 .
  \]
  Por tanto,
  \begin{align*}
    SS_{XY} &= \sumi (X_i-\media{X})(Y_i-\media{Y}) \\
            &= \sumi (X_i-\media{X})Y_i - \media{Y}\sumi (X_i-\media{X}) \\
            &= \sumi (X_i-\media{X})Y_i ,
  \end{align*}
  y dividiendo por $\SSX$,
  \begin{align*}
    \est{\beta}_1 &= \frac{SS_{XY}}{\SSX}
      = \sumi \frac{X_i-\media{X}}{\SSX}\,Y_i \\
    &= \sumi k_iY_i .
  \end{align*}

  \emph{Propiedad (i).}
  \[
    \sumi k_i = \frac{1}{\SSX}\sumi (X_i-\media{X}) = 0 .
  \]

  \emph{Propiedad (ii).} Restando el término nulo $\media{X}\displaystyle\sumi (X_i-\media{X})$,
  \begin{align*}
    \sumi k_iX_i &= \frac{1}{\SSX}\sumi (X_i-\media{X})X_i \\
    &= \frac{1}{\SSX}\left(\sumi (X_i-\media{X})X_i - \media{X}\sumi (X_i-\media{X})\right) \\
    &= \frac{1}{\SSX}\sumi (X_i-\media{X})(X_i-\media{X}) \\
    &= \frac{\SSX}{\SSX} = 1 .
  \end{align*}

  \emph{Propiedad (iii).}
  \begin{align*}
    \sumi k_i^2 &= \sumi \frac{(X_i-\media{X})^2}{\SSX^2} \\
    &= \frac{1}{\SSX^2}\sumi (X_i-\media{X})^2 \\
    &= \frac{\SSX}{\SSX^2} = \frac{1}{\SSX} .
  \end{align*}
\end{demostracion}

\begin{teorema}[Distribución de $\est{\beta}_1$][teo:t3-dist-b1]
  \[
    \est{\beta}_1 \sim \Normal\!\left(\beta_1,\ \frac{\sigma^2}{\SSX}\right),
    \qquad
    \frac{\est{\beta}_1-\beta_1}{\sqrt{MSE/\SSX}} \sim \tdist{n-2} .
  \]
\end{teorema}

\begin{demostracion}
  \emph{Normalidad.} Por la \cref{prop:t3-ki}, $\est{\beta}_1$ es combinación lineal de las
  $Y_i$, que son normales independientes; luego es normal.

  \emph{Insesgadez.} Usando las propiedades (i) y (ii) de los $k_i$,
  \begin{align*}
    \E(\est{\beta}_1) &= \E\left(\sumi k_iY_i\right) = \sumi k_i\E(Y_i) \\
    &= \sumi k_i(\beta_0+\beta_1X_i) \\
    &= \beta_0\sumi k_i + \beta_1\sumi k_iX_i \\
    &= \beta_0\cdot0+\beta_1\cdot1 = \beta_1 .
  \end{align*}

  \emph{Varianza.} Las $Y_i$ son independientes con varianza $\sigma^2$ y los $k_i$ son
  constantes, así que las covarianzas se anulan; usando la propiedad (iii),
  \begin{align*}
    \Var(\est{\beta}_1) &= \Var\left(\sumi k_iY_i\right) = \sumi k_i^2\Var(Y_i) \\
    &= \sigma^2\sumi k_i^2 \\
    &= \frac{\sigma^2}{\SSX} .
  \end{align*}
  Al estimar $\sigma^2$ por $MSE$, insesgado (\cref{prop:t2-mse}), se obtiene el estimador
  insesgado $\widehat{\Var}(\est{\beta}_1)=MSE/\SSX$.

  \emph{Distribución $t$.} Se utiliza que
  \[
    \frac{(n-2)MSE}{\sigma^2} = \frac{SSE}{\sigma^2} \sim \chisq{n-2}
  \]
  y que $SSE$ es independiente de $\est{\beta}_0$ y $\est{\beta}_1$ (resultado que se admite sin
  demostración). Entonces
  \begin{align*}
    \frac{\est{\beta}_1-\beta_1}{\sqrt{MSE/\SSX}}
    &= \frac{\dfrac{\est{\beta}_1-\beta_1}{\sqrt{\sigma^2/\SSX}}}
            {\sqrt{\dfrac{MSE}{\sigma^2}}}
     = \frac{\dfrac{\est{\beta}_1-\beta_1}{\sqrt{\sigma^2/\SSX}}}
            {\sqrt{\dfrac{SSE/\sigma^2}{n-2}}} ,
  \end{align*}
  que es el cociente entre una $\Normal(0,1)$ y la raíz de una $\chisq{n-2}$ independiente
  dividida por sus grados de libertad; luego sigue una $\tdist{n-2}$.
\end{demostracion}

\subsection{Intervalo de confianza para \texorpdfstring{$\beta_1$}{beta1}}

Del \cref{teo:t3-dist-b1},
\[
  P\left(\tq{n-2}{\alpha/2} \le \frac{\est{\beta}_1-\beta_1}{\sqrt{MSE/\SSX}}
  \le \tq{n-2}{1-\alpha/2}\right) = 1-\alpha .
\]
Por la simetría de la $t$ de Student, $\tq{n-2}{\alpha/2}=-\tq{n-2}{1-\alpha/2}$; despejando
$\beta_1$,
\[
  1-\alpha = P\left(\est{\beta}_1-\tq{n-2}{1-\alpha/2}\sqrt{\frac{MSE}{\SSX}}
  \le \beta_1 \le
  \est{\beta}_1+\tq{n-2}{1-\alpha/2}\sqrt{\frac{MSE}{\SSX}}\right).
\]

\begin{resumen}[Intervalo de confianza de nivel $1-\alpha$ para $\beta_1$]
  \[
    \beta_1 \in \left(\est{\beta}_1 \pm \tq{n-2}{1-\alpha/2}\sqrt{\frac{MSE}{\SSX}}\right).
  \]
  La cantidad $\sqrt{MSE/\SSX}$ es el \textbf{error estándar} de $\est{\beta}_1$.
\end{resumen}

\begin{ejemplo}[Volumen de madera][ej:arboles]
  Se quiere estimar la cosecha potencial de madera de un bosque. En una muestra de árboles se
  toman medidas no destructivas: HT (altura, en pies), DBH (diámetro del tronco a la altura del
  pecho, en pulgadas) y D16 (diámetro del tronco a 16 pies de altura, en pulgadas), con el fin
  de predecir el volumen de madera VOL (pies cúbicos) de cada árbol. Con $n=20$ árboles se
  ajusta un modelo de regresión lineal simple con respuesta VOL y regresor DBH, por ser la
  medida más fácil de obtener:
  \[
    \est{Y} = -45.57 + 6.916\,X, \qquad MSE = 36.587, \qquad \SSX = 64.512 .
  \]
  \begin{center}
    \begin{tikzpicture}
      \begin{axis}[width=7cm, height=5cm, xlabel={DBH (pulgadas)}, ylabel={VOL (pies$^3$)},
          xmin=9.5, xmax=20, ymin=15, ymax=105, label style={font=\small},
          tick label style={font=\footnotesize}]
        \addplot[puntos, mark size=1.3pt] table[x=dbh, y=vol] {datos/tema3/arboles.dat};
        \addplot[recta, domain=10:19.3] {-45.57+6.916*x};
      \end{axis}
    \end{tikzpicture}\\
    {\footnotesize (Datos reconstruidos a partir del gráfico de las diapositivas.)}
  \end{center}
  \textbf{IC del 95\,\% para $\beta_1$.} Con $\tq{18}{0.975}=2.101$ y error estándar
  $\sqrt{36.587/64.512}=0.7531$,
  \[
    \beta_1 \in (6.916 \pm 2.101\cdot0.7531) = (6.916\pm1.582) = (5.334,\ 8.498).
  \]
  Con una confianza del 95\,\%, el incremento medio poblacional del volumen de madera por cada
  pulgada de aumento del diámetro a la altura del pecho está entre 5.334 y 8.498 pies cúbicos.
  Como $0\notin(5.334,\,8.498)$, se rechaza $H_0\colon\beta_1=0$ a nivel $\alpha=0.05$: existe
  relación lineal entre DBH y VOL.
\end{ejemplo}

\subsection{Contrastes de hipótesis para \texorpdfstring{$\beta_1$}{beta1}}

Todos se basan en el estadístico del \cref{teo:t3-dist-b1}.

\paragraph{Contraste bilateral.} El más relevante es $H_0\colon\beta_1=0$ frente a
$H_1\colon\beta_1\neq0$, ya que si no se rechaza $H_0$ se concluye que no hay relación lineal
entre $X$ e $Y$. Se calcula el valor observado del estadístico
\[
  t^* = \frac{\est{\beta}_1-0}{\sqrt{MSE/\SSX}}
\]
y se compara con el cuantil de la $\tdist{n-2}$:
\begin{itemize}
  \item si $|t^*|\le\tq{n-2}{1-\alpha/2}$, no se rechaza $H_0$ a nivel $\alpha$;
  \item si $|t^*|>\tq{n-2}{1-\alpha/2}$, se rechaza $H_0$ a nivel $\alpha$.
\end{itemize}

\paragraph{Contraste unilateral.} Para $H_0\colon\beta_1\le0$ frente a $H_1\colon\beta_1>0$ se usa
el mismo estadístico, comparándolo ahora con el cuantil $1-\alpha$:
\begin{itemize}
  \item si $t^*\le\tq{n-2}{1-\alpha}$, no se rechaza $H_0$ a nivel $\alpha$;
  \item si $t^*>\tq{n-2}{1-\alpha}$, se rechaza $H_0$ a nivel $\alpha$.
\end{itemize}
(Análogamente para $H_1\colon\beta_1<0$, rechazando si $t^*<-\tq{n-2}{1-\alpha}$.)

\paragraph{Valor específico $b_1$.} Para $H_0\colon\beta_1=b_1$ frente a $H_1\colon\beta_1\neq b_1$
(o las versiones unilaterales) se usa
\[
  t^* = \frac{\est{\beta}_1-b_1}{\sqrt{MSE/\SSX}},
\]
con las mismas reglas de decisión.

\begin{ejemplo}[Volumen de madera: contrastes sobre $\beta_1$]
  \begin{itemize}
    \item \emph{Bilateral}, $H_0\colon\beta_1=0$:
      $t^* = 6.916/0.7531 = 9.18$. Como $|t^*|>\tq{18}{0.975}=2.101$, se rechaza $H_0$: hay
      relación lineal entre el diámetro del tronco y el volumen de madera.
    \item \emph{Unilateral}, $H_0\colon\beta_1\le0$ frente a $H_1\colon\beta_1>0$: el estadístico
      es el mismo, $t^*=9.18>\tq{18}{0.95}=1.73$, y se rechaza $H_0$: $\beta_1$ es positivo, es
      decir, el volumen aumenta significativamente con el diámetro.
    \item \emph{Valor específico}, $H_0\colon\beta_1=6$ frente a $H_1\colon\beta_1\neq6$:
      $t^* = (6.916-6)/0.7531 = 1.22 \le 2.101$, y no se rechaza $H_0$: no hay evidencia de que
      el incremento medio de volumen por pulgada de diámetro sea distinto de 6 pies cúbicos.
  \end{itemize}
\end{ejemplo}

\subsection{Inferencias sobre \texorpdfstring{$\beta_0$}{beta0}}

\begin{teorema}[Distribución de $\est{\beta}_0$][teo:t3-dist-b0]
  El estimador $\est{\beta}_0=\media{Y}-\est{\beta}_1\media{X}$ verifica
  \[
    \est{\beta}_0 \sim \Normal\!\left(\beta_0,\ \sigma^2\left(\frac1n+\frac{\media{X}^2}{\SSX}\right)\right),
    \qquad
    \frac{\est{\beta}_0-\beta_0}{\raizMSE{\frac1n+\frac{\media{X}^2}{\SSX}}} \sim \tdist{n-2}.
  \]
\end{teorema}

\begin{demostracion}
  \emph{Normalidad.} Escribiendo $\media{Y}$ y $\est{\beta}_1$ como combinaciones lineales de
  las $Y_i$ (\cref{prop:t3-ki}),
  \begin{align*}
    \est{\beta}_0 &= \media{Y}-\est{\beta}_1\media{X}
      = \sumi \frac1n\,Y_i - \media{X}\sumi k_iY_i \\
    &= \sumi\left(\frac1n - \media{X}k_i\right)Y_i \\
    &= \sumi\left(\frac1n - \frac{\media{X}(X_i-\media{X})}{\SSX}\right)Y_i ,
  \end{align*}
  que es combinación lineal de normales independientes; luego es normal.

  \emph{Insesgadez.} Como $\E(\est{\beta}_1)=\beta_1$,
  \begin{align*}
    \E(\est{\beta}_0) &= \E(\media{Y}) - \media{X}\,\E(\est{\beta}_1) \\
    &= \frac1n\sumi\E(Y_i) - \beta_1\media{X} \\
    &= \frac1n\sumi(\beta_0+\beta_1X_i) - \beta_1\media{X} \\
    &= \beta_0 + \beta_1\media{X} - \beta_1\media{X} \\
    &= \beta_0 .
  \end{align*}

  \emph{Varianza.} Por la independencia de las $Y_i$ y usando las propiedades (i) y (iii) de
  la \cref{prop:t3-ki},
  \begin{align*}
    \Var(\est{\beta}_0) &= \sumi\left(\frac1n-\media{X}k_i\right)^2\Var(Y_i) \\
    &= \sigma^2\sumi\left(\frac{1}{n^2} - \frac{2\media{X}}{n}\,k_i + \media{X}^2k_i^2\right) \\
    &= \sigma^2\left(\frac{n}{n^2} - \frac{2\media{X}}{n}\sumi k_i
       + \media{X}^2\sumi k_i^2\right) \\
    &= \sigma^2\left(\frac1n - 0 + \frac{\media{X}^2}{\SSX}\right)
     = \sigma^2\left(\frac1n + \frac{\media{X}^2}{\SSX}\right).
  \end{align*}
  Sustituyendo $\sigma^2$ por $MSE$ se obtiene el estimador insesgado
  \[
    \widehat{\Var}(\est{\beta}_0)=MSE\left(\frac1n+\frac{\media{X}^2}{\SSX}\right)
  \]
  y, por el mismo argumento que en el \cref{teo:t3-dist-b1}, la distribución $\tdist{n-2}$ del
  estadístico tipificado.
\end{demostracion}

\begin{resumen}[Inferencias sobre $\beta_0$]
  \begin{itemize}
    \item IC de nivel $1-\alpha$:
      $\displaystyle \beta_0 \in \left(\est{\beta}_0 \pm \tq{n-2}{1-\alpha/2}
      \raizMSE{\frac1n+\frac{\media{X}^2}{\SSX}}\right)$.
    \item Contraste de $H_0\colon\beta_0=b_0$:
      $\displaystyle t^* = \frac{\est{\beta}_0-b_0}{\raizMSE{\frac1n+\frac{\media{X}^2}{\SSX}}}$.
      En el contraste bilateral se rechaza $H_0$ si $|t^*|>\tq{n-2}{1-\alpha/2}$; en el
      unilateral ($H_1\colon\beta_0>b_0$), si $t^*>\tq{n-2}{1-\alpha}$.
  \end{itemize}
\end{resumen}

\subsection{Consideraciones sobre las inferencias}

\begin{itemize}
  \item \textbf{Normalidad.} Aunque las $Y_i$ no sean normales, las distribuciones de
    $\est{\beta}_0$ y $\est{\beta}_1$ son aproximadamente normales, en las condiciones del
    teorema central del límite.
  \item Se asume que $X$ es constante, y las inferencias son válidas en el \textbf{rango de
    valores de $X$ observado} en la muestra.
  \item Las varianzas de $\est{\beta}_0$ y $\est{\beta}_1$ dependen de $\SSX=\sum(X_i-\media{X})^2$:
    \textbf{mayor dispersión en $X$ implica menor variabilidad de los estimadores}. Esto es útil
    al planificar la recogida de datos.
\end{itemize}

\subsection{Inferencia simultánea: el problema de las comparaciones múltiples}
\label{sec:t3-bonferroni}

Hasta ahora se ha hecho inferencia sobre cada coeficiente por separado, fijando un nivel de
significación (por ejemplo, $\alpha=0.05$) en cada contraste o intervalo. Ese nivel está
garantizado para cada coeficiente por separado, pero no para el conjunto de las inferencias:
aunque fueran independientes, la confianza conjunta de dos intervalos del 95\,\% sería
$0.95^2\approx0.90$; y como $\est{\beta}_0$ y $\est{\beta}_1$ no son independientes (salvo si
$\media{X}=0$, \cref{prop:t5-var-simple}), ni siquiera puede calcularse como un producto.
Aparece así el problema de las \textbf{comparaciones múltiples}:
\begin{itemize}
  \item la confianza conjunta de varios IC simultáneos es menor que la de cada uno (menor del
    95\,\%);
  \item el nivel de significación del contraste simultáneo es mayor que el de cada contraste
    (mayor que 0.05): aumentan los falsos positivos.
\end{itemize}

\begin{ejemplo}[Probabilidad de falso positivo]
  Al lanzar un par de dados, llamemos ``error de tipo I'' a obtener dos unos, con probabilidad
  $\alpha=\frac16\cdot\frac16\approx0.028$. En una sola tirada la probabilidad de ese resultado
  es $\alpha$,
  pero al repetir las tiradas es cada vez más probable que alguna dé dos unos.

  Del mismo modo, si se realizan 20 contrastes independientes, cada uno a nivel $\alpha=0.05$,
  \begin{align*}
    P(\text{al menos un contraste significativo por azar})
    &= 1 - P(\text{ninguno significativo por azar}) \\
    &= 1-(1-0.05)^{20} \approx 0.64 .
  \end{align*}
\end{ejemplo}

Una solución a este problema es el \textbf{ajuste de Bonferroni}: modificar el nivel de
significación o de confianza de cada inferencia individual para conseguir el nivel deseado en la
inferencia simultánea.

\begin{proposicion}[Desigualdad de Bonferroni][prop:bonferroni]
  Para sucesos $A_1,\ldots,A_m$,
  \(
    P\left(\bigcup_{i=1}^m A_i\right) \le \sum_{i=1}^m P(A_i).
  \)
\end{proposicion}

\paragraph{Caso de $\beta_0$ y $\beta_1$.} Sea $A_0$ ($A_1$) el suceso ``el IC individual de
nivel $1-\alpha$ para $\beta_0$ ($\beta_1$) no contiene a $\beta_0$ ($\beta_1$)'', de modo que
$P(A_0)=P(A_1)=\alpha$. La probabilidad de que \emph{ambos} intervalos contengan a su parámetro
es
\[
  P(\overline{A_0}\cap\overline{A_1}) = 1-P(A_0\cup A_1)
  = 1-P(A_0)-P(A_1)+P(A_0\cap A_1) \ \ge\ 1-2\alpha ,
\]
porque $P(A_0\cap A_1)\ge0$. Así, dos IC individuales del 95\,\% garantizan una confianza
simultánea de al menos el 90\,\%, y dos IC individuales del 97.5\,\% garantizan una confianza
simultánea de al menos el 95\,\%.

\paragraph{Caso general.} Se construyen $m$ intervalos individuales de nivel $1-\alpha^*$,
\[
  C_i\colon\quad \beta_i \in \left(\est{\beta}_i \pm \tq{n-2}{1-\alpha^*/2}
  \sqrt{\widehat{\Var}(\est{\beta}_i)}\right),
  \qquad P(\beta_i\in C_i)=1-\alpha^* .
\]
Por la desigualdad de Bonferroni, la confianza simultánea cumple
\[
  P\left(\bigcap_{i=1}^m C_i\right)
  = 1-P\left(\bigcup_{i=1}^m \overline{C_i}\right)
  \ \ge\ 1-\sum_{i=1}^m P(\overline{C_i}) = 1-m\alpha^* .
\]

\begin{resumen}[Ajuste de Bonferroni]
  Para garantizar un nivel de confianza simultáneo $1-\alpha$ en $m$ intervalos (o un nivel de
  significación global $\alpha$ en $m$ contrastes) basta tomar en cada uno
  \[
    \alpha^* = \frac{\alpha}{m}.
  \]
\end{resumen}
